\teachcover{Lagrange, Galois,\\and reduced surds}{1}{
Which continued fractions repeat?
\[
 x>1,\qquad -1<x'<0
\]
}{
Begin with a quick reminder of the end of Lecture 1:
a repeating tail is a fixed point, and attaching a prefix
conjugates its matrix. Today's new work proves the converse
and the exact pure-periodicity criterion. Students need the
complete-quotient algorithm and convergence of $p_k/q_k$,
but no algebraic number theory. Source: book 3.8--3.9.
}

\teachslide{Recall: the complete quotient}{2}{
\[
 x_0=x,\qquad a_n=\floor{x_n},\qquad
 \key{x_{n+1}=\frac1{x_n-a_n}}
\]
\spacer
The number $x_n=[a_n;a_{n+1},\ldots]$ is the remaining tail.
\[
 x=[a_0;a_1,\ldots,a_{n-1},x_n].
\]
For $n\ge1$,
\[
 x_n>1,\qquad a_n\ge1.
\]
}{
Rebuild the notation rather than assuming everyone remembers
the subscript convention. The integer part of the original
number can be zero or negative. Every later complete quotient
is greater than 1, since it is the reciprocal of a nonzero
fractional part. A repeated complete quotient makes all future
digits repeat because the algorithm is deterministic.
Source: book 3.5 and 3.8.
}

\teachslide{The target theorem}{2}{
\textbf{Lagrange: the direction still to prove}
\[
 \key{\text{quadratic irrational}\quad\Longrightarrow\quad
 \text{eventually periodic}}
\]
\spacer
\textbf{Galois}
\[
 \key{\text{purely periodic}\quad\Longleftrightarrow\quad
 x>1,\ -1<x'<0}
\]
The second statement concerns quadratic irrationals.
\spacer
We will use integer states and a reversible step.
}{
Recall that $x'$ means the other root of the rational
quadratic equation, or the value obtained by reversing
the sign of its square root. The reducedness conditions
are not a conclusion about arbitrary real numbers;
they are defined for quadratic irrationals.
The key logical distinction is eventual versus pure
periodicity. Source: book Theorems 3.13 and 3.17.
External cross-check: Corless, \href{https://www.cecm.sfu.ca/organics/papers/corless/confrac/html/node6-an.shtml}{Periodic and Fixed Points of the Gauss Map}.
}

\teachslide{Conjugates follow the same algebra}{2}{
\[
 x_{n+1}=\frac1{x_n-a_n}
 \qquad\Longrightarrow\qquad
 \key{x_{n+1}'=\frac1{x_n'-a_n}}
\]
\spacer
The digit is still \key{$a_n=\floor{x_n}$}.

We do not take the integer part of $x_n'$.
\spacer
If $x_n'<0$ and $a_n\ge1$, then
\[
 x_n'-a_n<-1,\qquad -1<x_{n+1}'<0.
\]
}{
This is a common source of confusion. Conjugation changes
the algebraic number but leaves the rational integer $a_n$
fixed. It does not commute with taking a floor.
The inequality uses the floor of the positive complete
quotient. Ask students why the reciprocal of a number
less than $-1$ lies in $(-1,0)$. Once a negative
conjugate appears, the next tail is reduced.
Source: book 3.8, Step 1.
}

\teachslide{Reduced tails stay reduced}{2}{
Suppose
\[
 x_n>1,\qquad -1<x_n'<0.
\]
Then
\[
 x_{n+1}>1,\qquad -1<x_{n+1}'<0.
\]
\spacer
\key{Once reduced, every later complete quotient is reduced.}
\spacer
We still have to show that an arbitrary quadratic irrational
eventually reaches this region.
}{
The first inequality follows from the fractional-part
algorithm. The second follows from the previous slide.
Do not skip the last sentence: forward invariance is not
the same as eventual entry. This logical separation also
prevents a circular proof of Lagrange's theorem.
Source: book 3.8, Step 1.
}

\teachslide{Every quadratic irrational reaches a reduced tail}{3}{
The inverse prefix matrix gives, for $n\ge2$,
\[
 \key{x_n'=-\frac{q_{n-2}x'-p_{n-2}}
                    {q_{n-1}x'-p_{n-1}}.}
\]
The convergents tend to $x$, and $x'\ne x$.

For large $n$, numerator and denominator have the same sign.
\[
 x_n'<0\quad\text{eventually}.
\]
One more step puts the conjugate in $(-1,0)$.
}{
Give the full sign argument aloud. Since $p_k/q_k\to x$
and $x'\ne x$, the difference $x'-p_k/q_k$ eventually
has the fixed sign of $x'-x$. Multiplication by positive
$q_k$ preserves that sign. The two expressions in the
ratio therefore have the same sign, and the minus sign
makes $x_n'$ negative. Lecture 1's inverse formula gives
the displayed identity. A backup derives it in detail.
This covers negative original numbers too.
Source: book 3.8, Step 1.
}

\teachslide{An integer description of a surd}{2}{
Choose a fixed positive nonsquare integer $D$.
\[
 \key{x_n=\frac{P_n+\sqrt D}{Q_n}},
 \qquad P_n,Q_n\in\Z,\quad Q_n\ne0.
\]
We can choose the initial representation so that
\[
 \key{Q_n\mid D-P_n^2}
\]
holds at every step.
\spacer
$D$ need not be squarefree. Initially, $Q_n$ may be negative.
}{
Explain what the state remembers: the fixed square root
never changes, so only two integers need to evolve.
For completeness, if $Ax^2+Bx+C=0$, choose the common
sign of the coefficients so that
$x=(-B+\sqrt{B^2-4AC})/(2A)$.
Then $D=B^2-4AC$, $P_0=-B$, $Q_0=2A$, and
$D-P_0^2=-4AC$ is divisible by $Q_0$.
Do not automatically insist that the initial denominator
be positive. The first backup shows this normalization.
Source: book 3.8, Step 2.
}

\teachslide{The exact recurrence}{3}{
Set $a_n=\floor{x_n}$ and
\[
 \key{P_{n+1}=a_nQ_n-P_n.}
\]
Then
\[
 \begin{aligned}
 \frac1{x_n-a_n}
 &=\frac{Q_n}{\sqrt D-P_{n+1}}\\[5pt]
 &=\frac{P_{n+1}+\sqrt D}{(D-P_{n+1}^2)/Q_n}.
 \end{aligned}
\]
Thus
\[
 \key{Q_{n+1}=\frac{D-P_{n+1}^2}{Q_n}.}
\]
}{
Rationalize the denominator explicitly.
The sign in $P_{n+1}=a_nQ_n-P_n$ is easy to reverse
accidentally, so check it against
$x_n-a_n=(\sqrt D-P_{n+1})/Q_n$.
After rationalization, the numerator has precisely the
same form as before. We have not yet proved the new
denominator is an integer; that is the next slide.
Source: book 3.8, equation (3.4).
}

\teachslide{Why the recurrence stays integral}{2}{
\[
 P_{n+1}^2-P_n^2
 =a_nQ_n(a_nQ_n-2P_n).
\]
So $Q_n$ divides both $D-P_n^2$ and $D-P_{n+1}^2$.
\[
 \key{Q_{n+1}=(D-P_{n+1}^2)/Q_n\in\Z.}
\]
Also,
\[
 D-P_{n+1}^2=Q_nQ_{n+1}.
\]
This supplies the divisibility condition for the next step.
}{
This version uses divisibility directly, so students need
not already know congruence notation. Since $D$ is a
nonsquare integer, $D-P_{n+1}^2$ is never zero.
Thus $Q_{n+1}$ is a nonzero integer. The final equality
says that $Q_{n+1}$ itself divides
$D-P_{n+1}^2$, completing the induction.
Source: book 3.8, Step 2.
}

\teachslide{The bounds from reducedness}{3}{
For a reduced state,
\[
 x=\frac{P+\sqrt D}{Q}>1,\qquad
 -1<x'=\frac{P-\sqrt D}{Q}<0.
\]
Subtract:
\[
 x-x'=\frac{2\sqrt D}{Q}>1
 \quad\Longrightarrow\quad \key{0<Q<2\sqrt D.}
\]
The conjugate inequalities give
\[
 \sqrt D-Q<P<\sqrt D
 \quad\Longrightarrow\quad \key{|P|<\sqrt D.}
\]
}{
This is the compactness argument entirely in elementary
integer language. The positive difference $x-x'$ first
forces $Q>0$, so it is safe to multiply the conjugate
inequalities by $Q$. The lower bound becomes
$P>\sqrt D-Q>-\sqrt D$ because $Q<2\sqrt D$.
No uniform bound on the period length is needed.
Source: book 3.8, equation (3.5).
}

\teachslide{Finitely many states}{2}{
For fixed $D$, only finitely many integer pairs satisfy
\[
 \key{|P|<\sqrt D,\qquad 0<Q<2\sqrt D.}
\]
Every sufficiently late state lies in this finite set.

Each state determines its digit and its next state.
\[
 (P_i,Q_i)=(P_j,Q_j)
 \quad\Longrightarrow\quad
 a_{i+k}=a_{j+k}\quad(k\ge0).
\]
\key{This proves the converse direction of Lagrange's theorem.}
}{
Name the pigeonhole principle. Repetition of a pair means
repetition of the complete quotient, and deterministic
updating reproduces the entire future, not just one digit.
We have now proved both directions of Lagrange using
fixed points, convergence, rationalization, and finite
integer bounds. We still have not removed an initial
segment for a reduced starting state. Source: book 3.8.
}

\teachslide{The cycle for $\sqrt3$}{2}{
\begin{center}
\renewcommand{\arraystretch}{1.4}
\begin{tabular}{ccccc}
\toprule
$n$ & $P_n$ & $Q_n$ & $x_n$ & $a_n$\\
\midrule
$0$ & $0$ & $1$ & $\sqrt3$ & $1$\\
$1$ & $1$ & $2$ & $(1+\sqrt3)/2$ & $1$\\
$2$ & $1$ & $1$ & $1+\sqrt3$ & $2$\\
$3$ & $1$ & $2$ & $(1+\sqrt3)/2$ & $1$\\
\bottomrule
\end{tabular}
\end{center}
\[
 \key{\sqrt3=[1;\overline{1,2}]}
\]
The initial state lies outside the reduced cycle.
}{
Do one transition explicitly:
from $(P,Q,a)=(0,1,1)$ we get $P_{\rm new}=1$,
$Q_{\rm new}=2$. From $(1,2,1)$ we get $(1,1)$.
The equality of states at $n=1$ and $n=3$ certifies
period two. The finite table illustrates the proof,
but the bounds explain why every quadratic irrational
must eventually produce such a repetition.
Source: book Example 3.16.
}

\teachslide{Your turn: $\sqrt6$}{3}{
Start with $D=6$ and $(P_0,Q_0)=(0,1)$.
\[
 P_{n+1}=a_nQ_n-P_n,\qquad
 Q_{n+1}=\frac{6-P_{n+1}^2}{Q_n}.
\]
\spacer
Find the first repeated state.

Read off the continued fraction of $\sqrt6$.
\spacer
\key{Record the triples $(P_n,Q_n,a_n)$.}
}{
Give about two minutes of working time and one minute
to compare. The first digit is 2. The states are
$(0,1)$, $(2,2)$, $(2,1)$, $(2,2)$, with digits
$2,2,4,2$. If a student gets a negative denominator
after the first step, check the sign of $P_{n+1}$.
The interval $2<\sqrt6<3$ is enough to determine all
digits in this example without a decimal calculator.
This is a new worked classroom example using book (3.4).
}

\teachslide{The period of $\sqrt6$}{1}{
\begin{center}
\renewcommand{\arraystretch}{1.35}
\begin{tabular}{cccc}
\toprule
$n$ & $P_n$ & $Q_n$ & $a_n$\\
\midrule
$0$ & $0$ & $1$ & $2$\\
$1$ & $2$ & $2$ & $2$\\
$2$ & $2$ & $1$ & $4$\\
$3$ & $2$ & $2$ & $2$\\
\bottomrule
\end{tabular}
\end{center}
\[
 \key{\sqrt6=[2;\overline{2,4}]}
\]
}{
Keep the answer reveal brief. The first two reduced
complete quotients are $(2+\sqrt6)/2$ and $2+\sqrt6$.
Their conjugates are $(2-\sqrt6)/2$ and $2-\sqrt6$,
both in $(-1,0)$. This example gives a quick independent
check that the recurrence and reducedness conditions
are working together.
}

\teachslide{Finiteness allows an initial segment}{2}{
A function on the finite set $\{0,1\}$ can satisfy
\[
 F(0)=1,\qquad F(1)=1.
\]
Starting at 0 gives
\[
 0,\ 1,\ 1,\ 1,\ldots
\]
\spacer
Finiteness guarantees eventual repetition.

\key{To prove pure periodicity, we need to recover predecessors.}
}{
This is the conceptual turning point of the second lecture.
The reduced states form a finite forward-invariant set,
but that alone would not prove that an orbit starts on
its cycle. The little finite-map example exhibits the
missing possibility. Galois's criterion succeeds because
the continued-fraction step has unique reduced
predecessors. Source: book 3.9.
}

\teachslide{Galois's criterion}{2}{
For a quadratic irrational $x$,
\[
 \key{x\text{ is purely periodic}}
\]
\[
 \Longleftrightarrow
\]
\[
 \key{x>1\qquad\text{and}\qquad -1<x'<0.}
\]
\spacer
We first prove that a reduced starting point has no preperiod.
}{
Recall the definition from Lecture 1. We already know
that every reduced orbit stays reduced and eventually
repeats. The next two slides strengthen that conclusion
to pure periodicity. Afterward we prove necessity.
Keep the logical direction explicit throughout.
Source: book Theorem 3.17.
}

\teachslide{A successor remembers its reduced predecessor}{3}{
Suppose $u$ is reduced, $a=\floor u$, and
\[
 z=\frac1{u-a}.
\]
Conjugating gives
\[
 -\frac1{z'}=a-u',\qquad
 a<-\frac1{z'}<a+1.
\]
Therefore
\[
 \key{a=\floor{-1/z'},\qquad u=a+\frac1z.}
\]
Both the previous digit and the reduced predecessor are unique.
}{
The strict inequalities use precisely $-1<u'<0$.
This is why the definition of reducedness is so well
chosen. Stress that the assertion concerns reduced
predecessors. Without that restriction, many different
integers could be placed before the same tail.
For a quick check, take $z=(1+\sqrt3)/2$.
Then $-1/z'=1+\sqrt3$, so $a=2$ and the predecessor
is $1+\sqrt3$. Source: book 3.9, equation (3.6).
}

\teachslide{Rewinding a repeated state}{2}{
Suppose all $x_n$ are reduced and $x_i=x_j$ with $i<j$.
\[
 \begin{aligned}
 x_i&=x_j\\
 x_{i-1}&=x_{j-1}\\
 &\ \vdots\\
 x_0&=x_{j-i}.
 \end{aligned}
\]
Every backward step uses the unique reduced predecessor.
\spacer
\key{The complete-quotient sequence repeats from its start.}
}{
If $i=0$, the conclusion already holds. Otherwise apply
the predecessor formula repeatedly until reaching index
zero. Equal quadratic numbers have equal conjugates,
so the recovered digit and predecessor agree.
Now deterministic forward updating gives pure periodicity.
The point is not that we ran the ordinary floor algorithm
backward on arbitrary reals, but that reduced surds
carry enough information in their conjugates.
Source: book 3.9.
}

\teachslide{Why a pure period must be reduced}{2}{
Suppose the pure period has length $\ell$.

Lecture 1 shows that $x$ is quadratic.

All sufficiently late complete quotients are reduced.
\[
 \key{x_{k\ell}=x_0\qquad(k=1,2,\ldots)}
\]
Choose $k$ large enough.
\spacer
\key{The starting value is itself one of those reduced tails.}
}{
This is the other implication of Galois.
There is no circular appeal to Galois in the eventual-entry
argument: it used convergence of the ordinary convergents
and the inverse prefix matrix.
Together with the preceding two slides, this completes
the characterization. Source: book 3.9.
}

\teachslide{A reducedness test}{3}{
Which of these have purely periodic continued fractions?
\[
 \sqrt2,\qquad
 1+\sqrt3,\qquad
 \frac{1+\sqrt5}{2},\qquad
 \frac{\sqrt2}{2}
\]
\spacer
Check \key{both} inequalities:
\[
 x>1,\qquad -1<x'<0.
\]
\spacer
Which of the four are eventually periodic?
}{
Allow about 90 seconds, then ask for reasons rather than
just yes/no answers. The middle two are reduced.
For $\sqrt2$ the conjugate is below $-1$.
For $\sqrt2/2$ the original number is below 1.
All four are quadratic irrationals, so all four are
eventually periodic by Lagrange. Do not reveal the
next slide until the discussion finishes.
Source: book 3.9 and Exercise 23.
}

\teachslide{The test, checked}{2}{
\begin{center}
\renewcommand{\arraystretch}{1.4}
\begin{tabular}{lll}
\toprule
$x$ & $x'$ & Purely periodic?\\
\midrule
$\sqrt2$ & $-\sqrt2$ & No\\
$1+\sqrt3$ & $1-\sqrt3$ & Yes\\
$(1+\sqrt5)/2$ & $(1-\sqrt5)/2$ & Yes\\
$\sqrt2/2$ & $-\sqrt2/2$ & No\\
\bottomrule
\end{tabular}
\end{center}
\spacer
\key{All four are eventually periodic.}
}{
For the golden ratio, the inequalities follow from
$2<\sqrt5<3$. Its period consists of the digit 1.
For $1+\sqrt3$, the period is $2,1$.
Ask a student to distinguish the roles of the two theorems
in the final line. Source: book 3.9.
}

\teachslide{A visible prefix can belong to the period}{1}{
\[
 \key{[2;\overline{1,2}]=1+\sqrt3=[\overline{2,1}]}
\]
\spacer
Its conjugate is $1-\sqrt3\in(-1,0)$.

Galois detects pure periodicity independently of how
we first write the digit string.
}{
Return to the example from Lecture 1.
The earlier calculation established the value.
Galois now explains the classification.
A written bar need not mark the earliest possible
start of repetition. Source: book 3.9.
}

\teachslide{Lagrange and Galois}{2}{
\textbf{Lagrange}
\[
 \text{eventually periodic}
 \quad\Longleftrightarrow\quad
 \text{quadratic irrational}
\]
Fixed points prove one direction. Bounded integer states prove the other.
\spacer
\textbf{Galois}
\[
 \text{purely periodic}
 \quad\Longleftrightarrow\quad
 x>1,\ -1<x'<0
\]
Unique reduced predecessors remove the initial segment.
}{
Keep the quadratic-irrational hypothesis on the Galois
line explicit orally. The two lectures together used
only matrix algebra, convergence of convergents,
rationalization, inequalities, and the pigeonhole
principle. Mention optional links to Pell's equation
and reversal of the period only if time remains.
Both appear in the backup slides. Source: book 3.7--3.10.
}

\teachslide{Exit question}{1}{
In the proof, where did we use
\[
 \key{\text{finiteness?}}
\]
Where did we use
\[
 \key{\text{reversibility?}}
\]
\spacer
What can go wrong if only the first property holds?
}{
Expected answer: finiteness forces some repeated state
and therefore eventual periodicity. Unique reduced
predecessors let us rewind that repetition to index zero,
which gives pure periodicity. Without reversibility,
an orbit can have a transient initial segment, as in
$0,1,1,\ldots$. This is the 50-minute stopping point.
Suggested follow-up: book Exercises 24--26, or derive
the period of $\sqrt{13}$ using the exact recurrence.
}

\backups

\backupslide{Choosing the initial integer state}{
Start with $Ax^2+Bx+C=0$, where $A,B,C\in\Z$ and $A\ne0$.

Change all coefficient signs if needed so that
\[
 x=\frac{-B+\sqrt D}{2A},\qquad D=B^2-4AC>0.
\]
Then choose
\[
 P_0=-B,\qquad Q_0=2A.
\]
\[
 D-P_0^2=-4AC=Q_0(-2C).
\]
Thus $Q_0\mid D-P_0^2$. Its sign need not be positive.
}{
The discriminant is nonsquare because $x$ is irrational.
If the desired root initially has the minus sign before
the square root, replacing $(A,B,C)$ by $(-A,-B,-C)$
puts that root in the displayed plus-sign form.
For example, $x=-\sqrt2$ can be represented with
$D=8$, $P_0=0$, $Q_0=-2$.
This explains why a proof for all real quadratic
irrationals must allow initial negative denominators.
Source: book 3.8, Step 2.
}

\backupslide{The inverse-prefix formula}{
\[
 x=\frac{p_{n-1}x_n+p_{n-2}}
          {q_{n-1}x_n+q_{n-2}}.
\]
Solving for $x_n$ gives
\[
 x_n=-\frac{q_{n-2}x-p_{n-2}}
               {q_{n-1}x-p_{n-1}}.
\]
Conjugating gives
\[
 \key{x_n'=-\frac{q_{n-2}x'-p_{n-2}}
                     {q_{n-1}x'-p_{n-1}}.}
\]
For large $k$, $q_kx'-p_k=q_k(x'-p_k/q_k)$ has a fixed sign.
}{
Use this if the sign argument in main slide 6 needs
more time. The denominators are nonzero because the
conjugate is irrational and the convergents are rational.
For $n\ge2$, the relevant denominators $q_{n-2}$ and
$q_{n-1}$ are positive. Convergence toward $x$, not
any assumption of periodicity, gives the common sign.
Source: book 3.8, Step 1.
}

\backupslide{The conjugate reverses the period}{
If $x=[\overline{b_0,\ldots,b_{\ell-1}}]$, then
\[
 \key{-\frac1{x'}=[\overline{b_{\ell-1},\ldots,b_0}].}
\]
Each digit matrix is symmetric, so reversing a product
transposes it.
\[
 N=\mat ABC{D_0}:
 \quad Cx^2+(D_0-A)x-B=0
\]
Substitute $x'=-1/z$:
\[
 Bz^2+(D_0-A)z-C=0,
\]
the fixed-point equation for $N^T$.
}{
Both $B$ and $C$ are positive. The transformed quadratic
therefore has one positive root and one negative root.
By Galois, $-1/x'>0$; the reversed periodic continued
fraction is positive too, so they are the same root.
For a concrete check,
$[\overline{1,2,3}]=(4+\sqrt{37})/7$, while the reversed
period is $(4+\sqrt{37})/3$.
Source: book Corollary 3.19 and Exercise 26.
}

\backupslide{The same matrix produces Pell solutions}{
Lecture 1 gave
\[
 U=\mat2312,\qquad T_U(\sqrt3)=\sqrt3.
\]
Ordinary matrix multiplication gives
\[
 U\begin{pmatrix}p\\q\end{pmatrix}
   =\begin{pmatrix}2p+3q\\p+2q\end{pmatrix}.
\]
\[
 (2p+3q)^2-3(p+2q)^2=p^2-3q^2.
\]
Starting from $(1,0)$ produces
\[
 (2,1),\ (7,4),\ (26,15),\ldots
\]
Each satisfies $p^2-3q^2=1$.
}{
This is an optional connection, not an additional theorem
needed for either periodicity proof. The determinant-one
matrix acts both as a M\"obius transformation fixing
$\sqrt3$ and as an ordinary linear transformation
preserving the quadratic form $p^2-3q^2$.
Avoid claiming that this calculation alone proves the
classification of all Pell solutions.
Source: book 3.10 and Exercise 18.
}
